\documentclass[11pt,a4paper]{article}
% main (standalone preamble for editing)
% vim: nowrap: spell:
% ══════════════════════════════════════════════════════════════════════════════
% Speed Maths --- shared preamble
% Included by every sheet and answer file via:  % main (standalone preamble for editing)
% vim: nowrap: spell:
% ══════════════════════════════════════════════════════════════════════════════
% Speed Maths --- shared preamble
% Included by every sheet and answer file via:  % main (standalone preamble for editing)
% vim: nowrap: spell:
% ══════════════════════════════════════════════════════════════════════════════
% Speed Maths --- shared preamble
% Included by every sheet and answer file via:  \input{../../shared/preamble}
% Editing THIS file changes the look of every sheet/answer across every pillar.
% ══════════════════════════════════════════════════════════════════════════════

\usepackage[a4paper, top=25mm, bottom=25mm, left=22mm, right=22mm]{geometry}
\usepackage{amsmath, amssymb}
\usepackage{enumitem}
\usepackage{booktabs}
\usepackage{array}
\usepackage{parskip}
\usepackage{microtype}
\usepackage{fancyhdr}
\usepackage{titlesec}
\usepackage{mdframed}
\usepackage{xcolor}
\usepackage[hidelinks]{hyperref}
\usepackage{graphicx}
\usepackage{tikz}
\usepackage{pgfplots}
\pgfplotsset{compat=1.18}

% ── Colours ───────────────────────────────────────────────────────────────────
\definecolor{darkgray}{gray}{0.15}
\definecolor{medgray}{gray}{0.45}
\definecolor{lightgray}{gray}{0.93}
\definecolor{boxgray}{gray}{0.88}

% ── Section formatting ──────────────────────────────────────────────────────────
\titleformat{\section}[block]
  {\normalfont\large\bfseries}{}
  {0pt}{}[\vspace{2pt}\hrule\vspace{6pt}]
\titlespacing*{\section}{0pt}{14pt}{4pt}

% ── List formatting ─────────────────────────────────────────────────────────────
\setlist[enumerate,1]{label=\textbf{\arabic*.}, leftmargin=2em, itemsep=10pt}

% ── Branding constants (edit here, not per-file) ────────────────────────────────
\newcommand{\SpeedExamLine}{\textbf{Speed Maths:} Competition \& Exam Prep}
\newcommand{\SpeedCredit}{Series by \href{https://www.linkedin.com/in/cerealdev/}{David Akinyele-Aje}}

% ── Header / footer ──────────────────────────────────────────────────────────────
% Usage (once, after \input{../../shared/preamble} and before \begin{document}):
%   \SpeedHeader{Algebra}{3}
% First arg  = pillar name shown as "Speed <Pillar> --- Daily Drill"
% Second arg = worksheet number
\pagestyle{fancy}
\newcommand{\SpeedHeader}[2]{%
  \fancyhf{}%
  \renewcommand{\headrulewidth}{0.4pt}%
  \setlength{\headheight}{14pt}%
  \fancyhead[L]{\small\textbf{Speed #1 --- Daily Drill}}%
  \fancyhead[R]{\small Worksheet \##2}%
  \fancyfoot[L]{\small \SpeedExamLine}%
  \fancyfoot[C]{\small\thepage}%
  \fancyfoot[R]{\small \SpeedCredit}%
  \renewcommand{\footrulewidth}{0.4pt}%
}

% ── Title block (used at the top of every sheet AND answer file) ───────────────
% Usage: \SpeedTitleBlock{Daily Algebra Drill \#3}{Jane Doe}
% %% AGENTS: When generating a new sheet, ask the human contributor for the name they want credited in argument 2.
% For answer files, pass the "--- Answers \& Investigations" suffix in the arg.
\newcommand{\SpeedTitleBlock}[2]{%
  \begin{center}
    {\LARGE\bfseries #1}\\[4pt]
    {\normalsize\itshape Sheet by #2}\\[6pt]
    \hrule\vspace{4pt}
    \begin{tabular}{llcc}
      \toprule
      \textbf{Section} & \textbf{Topic} & \textbf{Questions} & \textbf{Target time}\\
      \midrule
      A & Rapid Recognition          & 10 & 2 min 30 sec \\
      B & Manipulation Drills        & 10 & 8 min        \\
      C & Substitution \& Structure  & 8  & 10 min       \\
      D & Challenge                  & 5  & 15 min       \\
      \bottomrule
    \end{tabular}
    \vspace{4pt}
    \hrule
  \end{center}
}

% ── Sheet metadata: topic + the day's new tools ────────────────────────────────
% Usage (immediately after \SpeedTitleBlock, sheets only):
%   \SpeedMeta{Expanding, factorising, and the identity toolkit}
%             {difference of squares, sum/product form, completing the square}
%
% Arg 1 = the sheet's topic in one line. The website lists this instead of
%         "Sheet 03", so write the thing a reader would actually search for.
% Arg 2 = comma-separated, at most five: the tools this sheet introduces that
%         earlier sheets in the pillar did not. These become the website's tags.
%         Leave out anything the reader already met — the tags are meant to say
%         what is new today, not to summarise the sheet.
%
% Both are parsed straight out of this file by tools/build_website.py, so the
% sheet stays the single source of truth and the site cannot drift from it.
%
% This renders NOTHING. It is a declaration for the build script, not a piece of
% the sheet's layout: adding it to a sheet must not change that sheet's PDF, so
% the 25 existing sheets can be annotated without a single one of them being
% reissued. If the toolkit should also appear on the page, that is a separate
% decision about the sheet design — make it deliberately, here, not by leaving
% this macro printing something nobody asked it to.
\newcommand{\SpeedMeta}[2]{}

% ── Answer / Method / Investigate-further boxes (answer files only) ────────────
\newmdenv[
  backgroundcolor=lightgray,
  linecolor=medgray,
  linewidth=0.5pt,
  leftmargin=0pt, rightmargin=0pt,
  innerleftmargin=8pt, innerrightmargin=8pt,
  innertopmargin=5pt, innerbottommargin=5pt,
  skipabove=4pt, skipbelow=4pt
]{ansbox}

\newmdenv[
  backgroundcolor=white,
  linecolor=darkgray,
  linewidth=0.8pt,
  leftline=true, rightline=false, topline=false, bottomline=false,
  leftmargin=0pt, rightmargin=0pt,
  innerleftmargin=8pt, innerrightmargin=6pt,
  innertopmargin=4pt, innerbottommargin=4pt,
  skipabove=4pt, skipbelow=6pt
]{invbox}

\newcommand{\ans}[1]{%
  \begin{ansbox}
    \textbf{Answer:} #1
  \end{ansbox}}

\newcommand{\method}[1]{%
  {\small\textbf{Method:} #1}\par}

\newcommand{\inv}[1]{%
  \begin{invbox}
    \textbf{\small Investigate further:} {\small #1}
  \end{invbox}}

% ── Misc helpers ────────────────────────────────────────────────────────────────
\newcommand{\qn}[1]{\textbf{#1.}\;}

% Mistake-linking: attach a Top 5 Patterns / Common Traps bullet to the question
% label(s) in this sheet where it actually shows up, e.g. \seealso{B6, D2}
\newcommand{\seealso}[1]{\hfill\textit{\footnotesize(see #1)}}

% ── Graph options macro ─────────────────────────────────────────────────────────
% Usage: \GraphOptions{tikz code for A}{tikz code for B}{tikz code for C}{tikz code for D}
\newcommand{\GraphOptions}[4]{%
  \begin{center}
    \begin{tabular}{cc}
      \textbf{A)} & \textbf{B)} \\
      #1 & #2 \\[10pt]
      \textbf{C)} & \textbf{D)} \\
      #3 & #4
    \end{tabular}
  \end{center}
}

% ── Closing motivational rule + cross-promo footer (sheets only) ────────────────
% Usage: \SpeedClosing{"Quote text."}
\newcommand{\SpeedClosing}[1]{%
  \vspace{20pt}
  \begin{center}
  \rule{0.6\textwidth}{0.4pt}\\[4pt]
  {\small\itshape #1}\\[4pt]
  \rule{0.6\textwidth}{0.4pt}
  \end{center}
  \begin{center}
  {\small Found a mistake or want to contribute a sheet?}\\
  {\small Visit the open-source project at \href{https://github.com/The-CerealDev/speed-maths}{github.com/The-CerealDev/speed-maths}}
  \end{center}
}

% Editing THIS file changes the look of every sheet/answer across every pillar.
% ══════════════════════════════════════════════════════════════════════════════

\usepackage[a4paper, top=25mm, bottom=25mm, left=22mm, right=22mm]{geometry}
\usepackage{amsmath, amssymb}
\usepackage{enumitem}
\usepackage{booktabs}
\usepackage{array}
\usepackage{parskip}
\usepackage{microtype}
\usepackage{fancyhdr}
\usepackage{titlesec}
\usepackage{mdframed}
\usepackage{xcolor}
\usepackage[hidelinks]{hyperref}
\usepackage{graphicx}
\usepackage{tikz}
\usepackage{pgfplots}
\pgfplotsset{compat=1.18}

% ── Colours ───────────────────────────────────────────────────────────────────
\definecolor{darkgray}{gray}{0.15}
\definecolor{medgray}{gray}{0.45}
\definecolor{lightgray}{gray}{0.93}
\definecolor{boxgray}{gray}{0.88}

% ── Section formatting ──────────────────────────────────────────────────────────
\titleformat{\section}[block]
  {\normalfont\large\bfseries}{}
  {0pt}{}[\vspace{2pt}\hrule\vspace{6pt}]
\titlespacing*{\section}{0pt}{14pt}{4pt}

% ── List formatting ─────────────────────────────────────────────────────────────
\setlist[enumerate,1]{label=\textbf{\arabic*.}, leftmargin=2em, itemsep=10pt}

% ── Branding constants (edit here, not per-file) ────────────────────────────────
\newcommand{\SpeedExamLine}{\textbf{Speed Maths:} Competition \& Exam Prep}
\newcommand{\SpeedCredit}{Series by \href{https://www.linkedin.com/in/cerealdev/}{David Akinyele-Aje}}

% ── Header / footer ──────────────────────────────────────────────────────────────
% Usage (once, after % main (standalone preamble for editing)
% vim: nowrap: spell:
% ══════════════════════════════════════════════════════════════════════════════
% Speed Maths --- shared preamble
% Included by every sheet and answer file via:  \input{../../shared/preamble}
% Editing THIS file changes the look of every sheet/answer across every pillar.
% ══════════════════════════════════════════════════════════════════════════════

\usepackage[a4paper, top=25mm, bottom=25mm, left=22mm, right=22mm]{geometry}
\usepackage{amsmath, amssymb}
\usepackage{enumitem}
\usepackage{booktabs}
\usepackage{array}
\usepackage{parskip}
\usepackage{microtype}
\usepackage{fancyhdr}
\usepackage{titlesec}
\usepackage{mdframed}
\usepackage{xcolor}
\usepackage[hidelinks]{hyperref}
\usepackage{graphicx}
\usepackage{tikz}
\usepackage{pgfplots}
\pgfplotsset{compat=1.18}

% ── Colours ───────────────────────────────────────────────────────────────────
\definecolor{darkgray}{gray}{0.15}
\definecolor{medgray}{gray}{0.45}
\definecolor{lightgray}{gray}{0.93}
\definecolor{boxgray}{gray}{0.88}

% ── Section formatting ──────────────────────────────────────────────────────────
\titleformat{\section}[block]
  {\normalfont\large\bfseries}{}
  {0pt}{}[\vspace{2pt}\hrule\vspace{6pt}]
\titlespacing*{\section}{0pt}{14pt}{4pt}

% ── List formatting ─────────────────────────────────────────────────────────────
\setlist[enumerate,1]{label=\textbf{\arabic*.}, leftmargin=2em, itemsep=10pt}

% ── Branding constants (edit here, not per-file) ────────────────────────────────
\newcommand{\SpeedExamLine}{\textbf{Speed Maths:} Competition \& Exam Prep}
\newcommand{\SpeedCredit}{Series by \href{https://www.linkedin.com/in/cerealdev/}{David Akinyele-Aje}}

% ── Header / footer ──────────────────────────────────────────────────────────────
% Usage (once, after \input{../../shared/preamble} and before \begin{document}):
%   \SpeedHeader{Algebra}{3}
% First arg  = pillar name shown as "Speed <Pillar> --- Daily Drill"
% Second arg = worksheet number
\pagestyle{fancy}
\newcommand{\SpeedHeader}[2]{%
  \fancyhf{}%
  \renewcommand{\headrulewidth}{0.4pt}%
  \setlength{\headheight}{14pt}%
  \fancyhead[L]{\small\textbf{Speed #1 --- Daily Drill}}%
  \fancyhead[R]{\small Worksheet \##2}%
  \fancyfoot[L]{\small \SpeedExamLine}%
  \fancyfoot[C]{\small\thepage}%
  \fancyfoot[R]{\small \SpeedCredit}%
  \renewcommand{\footrulewidth}{0.4pt}%
}

% ── Title block (used at the top of every sheet AND answer file) ───────────────
% Usage: \SpeedTitleBlock{Daily Algebra Drill \#3}{Jane Doe}
% %% AGENTS: When generating a new sheet, ask the human contributor for the name they want credited in argument 2.
% For answer files, pass the "--- Answers \& Investigations" suffix in the arg.
\newcommand{\SpeedTitleBlock}[2]{%
  \begin{center}
    {\LARGE\bfseries #1}\\[4pt]
    {\normalsize\itshape Sheet by #2}\\[6pt]
    \hrule\vspace{4pt}
    \begin{tabular}{llcc}
      \toprule
      \textbf{Section} & \textbf{Topic} & \textbf{Questions} & \textbf{Target time}\\
      \midrule
      A & Rapid Recognition          & 10 & 2 min 30 sec \\
      B & Manipulation Drills        & 10 & 8 min        \\
      C & Substitution \& Structure  & 8  & 10 min       \\
      D & Challenge                  & 5  & 15 min       \\
      \bottomrule
    \end{tabular}
    \vspace{4pt}
    \hrule
  \end{center}
}

% ── Sheet metadata: topic + the day's new tools ────────────────────────────────
% Usage (immediately after \SpeedTitleBlock, sheets only):
%   \SpeedMeta{Expanding, factorising, and the identity toolkit}
%             {difference of squares, sum/product form, completing the square}
%
% Arg 1 = the sheet's topic in one line. The website lists this instead of
%         "Sheet 03", so write the thing a reader would actually search for.
% Arg 2 = comma-separated, at most five: the tools this sheet introduces that
%         earlier sheets in the pillar did not. These become the website's tags.
%         Leave out anything the reader already met — the tags are meant to say
%         what is new today, not to summarise the sheet.
%
% Both are parsed straight out of this file by tools/build_website.py, so the
% sheet stays the single source of truth and the site cannot drift from it.
%
% This renders NOTHING. It is a declaration for the build script, not a piece of
% the sheet's layout: adding it to a sheet must not change that sheet's PDF, so
% the 25 existing sheets can be annotated without a single one of them being
% reissued. If the toolkit should also appear on the page, that is a separate
% decision about the sheet design — make it deliberately, here, not by leaving
% this macro printing something nobody asked it to.
\newcommand{\SpeedMeta}[2]{}

% ── Answer / Method / Investigate-further boxes (answer files only) ────────────
\newmdenv[
  backgroundcolor=lightgray,
  linecolor=medgray,
  linewidth=0.5pt,
  leftmargin=0pt, rightmargin=0pt,
  innerleftmargin=8pt, innerrightmargin=8pt,
  innertopmargin=5pt, innerbottommargin=5pt,
  skipabove=4pt, skipbelow=4pt
]{ansbox}

\newmdenv[
  backgroundcolor=white,
  linecolor=darkgray,
  linewidth=0.8pt,
  leftline=true, rightline=false, topline=false, bottomline=false,
  leftmargin=0pt, rightmargin=0pt,
  innerleftmargin=8pt, innerrightmargin=6pt,
  innertopmargin=4pt, innerbottommargin=4pt,
  skipabove=4pt, skipbelow=6pt
]{invbox}

\newcommand{\ans}[1]{%
  \begin{ansbox}
    \textbf{Answer:} #1
  \end{ansbox}}

\newcommand{\method}[1]{%
  {\small\textbf{Method:} #1}\par}

\newcommand{\inv}[1]{%
  \begin{invbox}
    \textbf{\small Investigate further:} {\small #1}
  \end{invbox}}

% ── Misc helpers ────────────────────────────────────────────────────────────────
\newcommand{\qn}[1]{\textbf{#1.}\;}

% Mistake-linking: attach a Top 5 Patterns / Common Traps bullet to the question
% label(s) in this sheet where it actually shows up, e.g. \seealso{B6, D2}
\newcommand{\seealso}[1]{\hfill\textit{\footnotesize(see #1)}}

% ── Graph options macro ─────────────────────────────────────────────────────────
% Usage: \GraphOptions{tikz code for A}{tikz code for B}{tikz code for C}{tikz code for D}
\newcommand{\GraphOptions}[4]{%
  \begin{center}
    \begin{tabular}{cc}
      \textbf{A)} & \textbf{B)} \\
      #1 & #2 \\[10pt]
      \textbf{C)} & \textbf{D)} \\
      #3 & #4
    \end{tabular}
  \end{center}
}

% ── Closing motivational rule + cross-promo footer (sheets only) ────────────────
% Usage: \SpeedClosing{"Quote text."}
\newcommand{\SpeedClosing}[1]{%
  \vspace{20pt}
  \begin{center}
  \rule{0.6\textwidth}{0.4pt}\\[4pt]
  {\small\itshape #1}\\[4pt]
  \rule{0.6\textwidth}{0.4pt}
  \end{center}
  \begin{center}
  {\small Found a mistake or want to contribute a sheet?}\\
  {\small Visit the open-source project at \href{https://github.com/The-CerealDev/speed-maths}{github.com/The-CerealDev/speed-maths}}
  \end{center}
}
 and before \begin{document}):
%   \SpeedHeader{Algebra}{3}
% First arg  = pillar name shown as "Speed <Pillar> --- Daily Drill"
% Second arg = worksheet number
\pagestyle{fancy}
\newcommand{\SpeedHeader}[2]{%
  \fancyhf{}%
  \renewcommand{\headrulewidth}{0.4pt}%
  \setlength{\headheight}{14pt}%
  \fancyhead[L]{\small\textbf{Speed #1 --- Daily Drill}}%
  \fancyhead[R]{\small Worksheet \##2}%
  \fancyfoot[L]{\small \SpeedExamLine}%
  \fancyfoot[C]{\small\thepage}%
  \fancyfoot[R]{\small \SpeedCredit}%
  \renewcommand{\footrulewidth}{0.4pt}%
}

% ── Title block (used at the top of every sheet AND answer file) ───────────────
% Usage: \SpeedTitleBlock{Daily Algebra Drill \#3}{Jane Doe}
% %% AGENTS: When generating a new sheet, ask the human contributor for the name they want credited in argument 2.
% For answer files, pass the "--- Answers \& Investigations" suffix in the arg.
\newcommand{\SpeedTitleBlock}[2]{%
  \begin{center}
    {\LARGE\bfseries #1}\\[4pt]
    {\normalsize\itshape Sheet by #2}\\[6pt]
    \hrule\vspace{4pt}
    \begin{tabular}{llcc}
      \toprule
      \textbf{Section} & \textbf{Topic} & \textbf{Questions} & \textbf{Target time}\\
      \midrule
      A & Rapid Recognition          & 10 & 2 min 30 sec \\
      B & Manipulation Drills        & 10 & 8 min        \\
      C & Substitution \& Structure  & 8  & 10 min       \\
      D & Challenge                  & 5  & 15 min       \\
      \bottomrule
    \end{tabular}
    \vspace{4pt}
    \hrule
  \end{center}
}

% ── Sheet metadata: topic + the day's new tools ────────────────────────────────
% Usage (immediately after \SpeedTitleBlock, sheets only):
%   \SpeedMeta{Expanding, factorising, and the identity toolkit}
%             {difference of squares, sum/product form, completing the square}
%
% Arg 1 = the sheet's topic in one line. The website lists this instead of
%         "Sheet 03", so write the thing a reader would actually search for.
% Arg 2 = comma-separated, at most five: the tools this sheet introduces that
%         earlier sheets in the pillar did not. These become the website's tags.
%         Leave out anything the reader already met — the tags are meant to say
%         what is new today, not to summarise the sheet.
%
% Both are parsed straight out of this file by tools/build_website.py, so the
% sheet stays the single source of truth and the site cannot drift from it.
%
% This renders NOTHING. It is a declaration for the build script, not a piece of
% the sheet's layout: adding it to a sheet must not change that sheet's PDF, so
% the 25 existing sheets can be annotated without a single one of them being
% reissued. If the toolkit should also appear on the page, that is a separate
% decision about the sheet design — make it deliberately, here, not by leaving
% this macro printing something nobody asked it to.
\newcommand{\SpeedMeta}[2]{}

% ── Answer / Method / Investigate-further boxes (answer files only) ────────────
\newmdenv[
  backgroundcolor=lightgray,
  linecolor=medgray,
  linewidth=0.5pt,
  leftmargin=0pt, rightmargin=0pt,
  innerleftmargin=8pt, innerrightmargin=8pt,
  innertopmargin=5pt, innerbottommargin=5pt,
  skipabove=4pt, skipbelow=4pt
]{ansbox}

\newmdenv[
  backgroundcolor=white,
  linecolor=darkgray,
  linewidth=0.8pt,
  leftline=true, rightline=false, topline=false, bottomline=false,
  leftmargin=0pt, rightmargin=0pt,
  innerleftmargin=8pt, innerrightmargin=6pt,
  innertopmargin=4pt, innerbottommargin=4pt,
  skipabove=4pt, skipbelow=6pt
]{invbox}

\newcommand{\ans}[1]{%
  \begin{ansbox}
    \textbf{Answer:} #1
  \end{ansbox}}

\newcommand{\method}[1]{%
  {\small\textbf{Method:} #1}\par}

\newcommand{\inv}[1]{%
  \begin{invbox}
    \textbf{\small Investigate further:} {\small #1}
  \end{invbox}}

% ── Misc helpers ────────────────────────────────────────────────────────────────
\newcommand{\qn}[1]{\textbf{#1.}\;}

% Mistake-linking: attach a Top 5 Patterns / Common Traps bullet to the question
% label(s) in this sheet where it actually shows up, e.g. \seealso{B6, D2}
\newcommand{\seealso}[1]{\hfill\textit{\footnotesize(see #1)}}

% ── Graph options macro ─────────────────────────────────────────────────────────
% Usage: \GraphOptions{tikz code for A}{tikz code for B}{tikz code for C}{tikz code for D}
\newcommand{\GraphOptions}[4]{%
  \begin{center}
    \begin{tabular}{cc}
      \textbf{A)} & \textbf{B)} \\
      #1 & #2 \\[10pt]
      \textbf{C)} & \textbf{D)} \\
      #3 & #4
    \end{tabular}
  \end{center}
}

% ── Closing motivational rule + cross-promo footer (sheets only) ────────────────
% Usage: \SpeedClosing{"Quote text."}
\newcommand{\SpeedClosing}[1]{%
  \vspace{20pt}
  \begin{center}
  \rule{0.6\textwidth}{0.4pt}\\[4pt]
  {\small\itshape #1}\\[4pt]
  \rule{0.6\textwidth}{0.4pt}
  \end{center}
  \begin{center}
  {\small Found a mistake or want to contribute a sheet?}\\
  {\small Visit the open-source project at \href{https://github.com/The-CerealDev/speed-maths}{github.com/The-CerealDev/speed-maths}}
  \end{center}
}

% Editing THIS file changes the look of every sheet/answer across every pillar.
% ══════════════════════════════════════════════════════════════════════════════

\usepackage[a4paper, top=25mm, bottom=25mm, left=22mm, right=22mm]{geometry}
\usepackage{amsmath, amssymb}
\usepackage{enumitem}
\usepackage{booktabs}
\usepackage{array}
\usepackage{parskip}
\usepackage{microtype}
\usepackage{fancyhdr}
\usepackage{titlesec}
\usepackage{mdframed}
\usepackage{xcolor}
\usepackage[hidelinks]{hyperref}
\usepackage{graphicx}
\usepackage{tikz}
\usepackage{pgfplots}
\pgfplotsset{compat=1.18}

% ── Colours ───────────────────────────────────────────────────────────────────
\definecolor{darkgray}{gray}{0.15}
\definecolor{medgray}{gray}{0.45}
\definecolor{lightgray}{gray}{0.93}
\definecolor{boxgray}{gray}{0.88}

% ── Section formatting ──────────────────────────────────────────────────────────
\titleformat{\section}[block]
  {\normalfont\large\bfseries}{}
  {0pt}{}[\vspace{2pt}\hrule\vspace{6pt}]
\titlespacing*{\section}{0pt}{14pt}{4pt}

% ── List formatting ─────────────────────────────────────────────────────────────
\setlist[enumerate,1]{label=\textbf{\arabic*.}, leftmargin=2em, itemsep=10pt}

% ── Branding constants (edit here, not per-file) ────────────────────────────────
\newcommand{\SpeedExamLine}{\textbf{Speed Maths:} Competition \& Exam Prep}
\newcommand{\SpeedCredit}{Series by \href{https://www.linkedin.com/in/cerealdev/}{David Akinyele-Aje}}

% ── Header / footer ──────────────────────────────────────────────────────────────
% Usage (once, after % main (standalone preamble for editing)
% vim: nowrap: spell:
% ══════════════════════════════════════════════════════════════════════════════
% Speed Maths --- shared preamble
% Included by every sheet and answer file via:  % main (standalone preamble for editing)
% vim: nowrap: spell:
% ══════════════════════════════════════════════════════════════════════════════
% Speed Maths --- shared preamble
% Included by every sheet and answer file via:  \input{../../shared/preamble}
% Editing THIS file changes the look of every sheet/answer across every pillar.
% ══════════════════════════════════════════════════════════════════════════════

\usepackage[a4paper, top=25mm, bottom=25mm, left=22mm, right=22mm]{geometry}
\usepackage{amsmath, amssymb}
\usepackage{enumitem}
\usepackage{booktabs}
\usepackage{array}
\usepackage{parskip}
\usepackage{microtype}
\usepackage{fancyhdr}
\usepackage{titlesec}
\usepackage{mdframed}
\usepackage{xcolor}
\usepackage[hidelinks]{hyperref}
\usepackage{graphicx}
\usepackage{tikz}
\usepackage{pgfplots}
\pgfplotsset{compat=1.18}

% ── Colours ───────────────────────────────────────────────────────────────────
\definecolor{darkgray}{gray}{0.15}
\definecolor{medgray}{gray}{0.45}
\definecolor{lightgray}{gray}{0.93}
\definecolor{boxgray}{gray}{0.88}

% ── Section formatting ──────────────────────────────────────────────────────────
\titleformat{\section}[block]
  {\normalfont\large\bfseries}{}
  {0pt}{}[\vspace{2pt}\hrule\vspace{6pt}]
\titlespacing*{\section}{0pt}{14pt}{4pt}

% ── List formatting ─────────────────────────────────────────────────────────────
\setlist[enumerate,1]{label=\textbf{\arabic*.}, leftmargin=2em, itemsep=10pt}

% ── Branding constants (edit here, not per-file) ────────────────────────────────
\newcommand{\SpeedExamLine}{\textbf{Speed Maths:} Competition \& Exam Prep}
\newcommand{\SpeedCredit}{Series by \href{https://www.linkedin.com/in/cerealdev/}{David Akinyele-Aje}}

% ── Header / footer ──────────────────────────────────────────────────────────────
% Usage (once, after \input{../../shared/preamble} and before \begin{document}):
%   \SpeedHeader{Algebra}{3}
% First arg  = pillar name shown as "Speed <Pillar> --- Daily Drill"
% Second arg = worksheet number
\pagestyle{fancy}
\newcommand{\SpeedHeader}[2]{%
  \fancyhf{}%
  \renewcommand{\headrulewidth}{0.4pt}%
  \setlength{\headheight}{14pt}%
  \fancyhead[L]{\small\textbf{Speed #1 --- Daily Drill}}%
  \fancyhead[R]{\small Worksheet \##2}%
  \fancyfoot[L]{\small \SpeedExamLine}%
  \fancyfoot[C]{\small\thepage}%
  \fancyfoot[R]{\small \SpeedCredit}%
  \renewcommand{\footrulewidth}{0.4pt}%
}

% ── Title block (used at the top of every sheet AND answer file) ───────────────
% Usage: \SpeedTitleBlock{Daily Algebra Drill \#3}{Jane Doe}
% %% AGENTS: When generating a new sheet, ask the human contributor for the name they want credited in argument 2.
% For answer files, pass the "--- Answers \& Investigations" suffix in the arg.
\newcommand{\SpeedTitleBlock}[2]{%
  \begin{center}
    {\LARGE\bfseries #1}\\[4pt]
    {\normalsize\itshape Sheet by #2}\\[6pt]
    \hrule\vspace{4pt}
    \begin{tabular}{llcc}
      \toprule
      \textbf{Section} & \textbf{Topic} & \textbf{Questions} & \textbf{Target time}\\
      \midrule
      A & Rapid Recognition          & 10 & 2 min 30 sec \\
      B & Manipulation Drills        & 10 & 8 min        \\
      C & Substitution \& Structure  & 8  & 10 min       \\
      D & Challenge                  & 5  & 15 min       \\
      \bottomrule
    \end{tabular}
    \vspace{4pt}
    \hrule
  \end{center}
}

% ── Sheet metadata: topic + the day's new tools ────────────────────────────────
% Usage (immediately after \SpeedTitleBlock, sheets only):
%   \SpeedMeta{Expanding, factorising, and the identity toolkit}
%             {difference of squares, sum/product form, completing the square}
%
% Arg 1 = the sheet's topic in one line. The website lists this instead of
%         "Sheet 03", so write the thing a reader would actually search for.
% Arg 2 = comma-separated, at most five: the tools this sheet introduces that
%         earlier sheets in the pillar did not. These become the website's tags.
%         Leave out anything the reader already met — the tags are meant to say
%         what is new today, not to summarise the sheet.
%
% Both are parsed straight out of this file by tools/build_website.py, so the
% sheet stays the single source of truth and the site cannot drift from it.
%
% This renders NOTHING. It is a declaration for the build script, not a piece of
% the sheet's layout: adding it to a sheet must not change that sheet's PDF, so
% the 25 existing sheets can be annotated without a single one of them being
% reissued. If the toolkit should also appear on the page, that is a separate
% decision about the sheet design — make it deliberately, here, not by leaving
% this macro printing something nobody asked it to.
\newcommand{\SpeedMeta}[2]{}

% ── Answer / Method / Investigate-further boxes (answer files only) ────────────
\newmdenv[
  backgroundcolor=lightgray,
  linecolor=medgray,
  linewidth=0.5pt,
  leftmargin=0pt, rightmargin=0pt,
  innerleftmargin=8pt, innerrightmargin=8pt,
  innertopmargin=5pt, innerbottommargin=5pt,
  skipabove=4pt, skipbelow=4pt
]{ansbox}

\newmdenv[
  backgroundcolor=white,
  linecolor=darkgray,
  linewidth=0.8pt,
  leftline=true, rightline=false, topline=false, bottomline=false,
  leftmargin=0pt, rightmargin=0pt,
  innerleftmargin=8pt, innerrightmargin=6pt,
  innertopmargin=4pt, innerbottommargin=4pt,
  skipabove=4pt, skipbelow=6pt
]{invbox}

\newcommand{\ans}[1]{%
  \begin{ansbox}
    \textbf{Answer:} #1
  \end{ansbox}}

\newcommand{\method}[1]{%
  {\small\textbf{Method:} #1}\par}

\newcommand{\inv}[1]{%
  \begin{invbox}
    \textbf{\small Investigate further:} {\small #1}
  \end{invbox}}

% ── Misc helpers ────────────────────────────────────────────────────────────────
\newcommand{\qn}[1]{\textbf{#1.}\;}

% Mistake-linking: attach a Top 5 Patterns / Common Traps bullet to the question
% label(s) in this sheet where it actually shows up, e.g. \seealso{B6, D2}
\newcommand{\seealso}[1]{\hfill\textit{\footnotesize(see #1)}}

% ── Graph options macro ─────────────────────────────────────────────────────────
% Usage: \GraphOptions{tikz code for A}{tikz code for B}{tikz code for C}{tikz code for D}
\newcommand{\GraphOptions}[4]{%
  \begin{center}
    \begin{tabular}{cc}
      \textbf{A)} & \textbf{B)} \\
      #1 & #2 \\[10pt]
      \textbf{C)} & \textbf{D)} \\
      #3 & #4
    \end{tabular}
  \end{center}
}

% ── Closing motivational rule + cross-promo footer (sheets only) ────────────────
% Usage: \SpeedClosing{"Quote text."}
\newcommand{\SpeedClosing}[1]{%
  \vspace{20pt}
  \begin{center}
  \rule{0.6\textwidth}{0.4pt}\\[4pt]
  {\small\itshape #1}\\[4pt]
  \rule{0.6\textwidth}{0.4pt}
  \end{center}
  \begin{center}
  {\small Found a mistake or want to contribute a sheet?}\\
  {\small Visit the open-source project at \href{https://github.com/The-CerealDev/speed-maths}{github.com/The-CerealDev/speed-maths}}
  \end{center}
}

% Editing THIS file changes the look of every sheet/answer across every pillar.
% ══════════════════════════════════════════════════════════════════════════════

\usepackage[a4paper, top=25mm, bottom=25mm, left=22mm, right=22mm]{geometry}
\usepackage{amsmath, amssymb}
\usepackage{enumitem}
\usepackage{booktabs}
\usepackage{array}
\usepackage{parskip}
\usepackage{microtype}
\usepackage{fancyhdr}
\usepackage{titlesec}
\usepackage{mdframed}
\usepackage{xcolor}
\usepackage[hidelinks]{hyperref}
\usepackage{graphicx}
\usepackage{tikz}
\usepackage{pgfplots}
\pgfplotsset{compat=1.18}

% ── Colours ───────────────────────────────────────────────────────────────────
\definecolor{darkgray}{gray}{0.15}
\definecolor{medgray}{gray}{0.45}
\definecolor{lightgray}{gray}{0.93}
\definecolor{boxgray}{gray}{0.88}

% ── Section formatting ──────────────────────────────────────────────────────────
\titleformat{\section}[block]
  {\normalfont\large\bfseries}{}
  {0pt}{}[\vspace{2pt}\hrule\vspace{6pt}]
\titlespacing*{\section}{0pt}{14pt}{4pt}

% ── List formatting ─────────────────────────────────────────────────────────────
\setlist[enumerate,1]{label=\textbf{\arabic*.}, leftmargin=2em, itemsep=10pt}

% ── Branding constants (edit here, not per-file) ────────────────────────────────
\newcommand{\SpeedExamLine}{\textbf{Speed Maths:} Competition \& Exam Prep}
\newcommand{\SpeedCredit}{Series by \href{https://www.linkedin.com/in/cerealdev/}{David Akinyele-Aje}}

% ── Header / footer ──────────────────────────────────────────────────────────────
% Usage (once, after % main (standalone preamble for editing)
% vim: nowrap: spell:
% ══════════════════════════════════════════════════════════════════════════════
% Speed Maths --- shared preamble
% Included by every sheet and answer file via:  \input{../../shared/preamble}
% Editing THIS file changes the look of every sheet/answer across every pillar.
% ══════════════════════════════════════════════════════════════════════════════

\usepackage[a4paper, top=25mm, bottom=25mm, left=22mm, right=22mm]{geometry}
\usepackage{amsmath, amssymb}
\usepackage{enumitem}
\usepackage{booktabs}
\usepackage{array}
\usepackage{parskip}
\usepackage{microtype}
\usepackage{fancyhdr}
\usepackage{titlesec}
\usepackage{mdframed}
\usepackage{xcolor}
\usepackage[hidelinks]{hyperref}
\usepackage{graphicx}
\usepackage{tikz}
\usepackage{pgfplots}
\pgfplotsset{compat=1.18}

% ── Colours ───────────────────────────────────────────────────────────────────
\definecolor{darkgray}{gray}{0.15}
\definecolor{medgray}{gray}{0.45}
\definecolor{lightgray}{gray}{0.93}
\definecolor{boxgray}{gray}{0.88}

% ── Section formatting ──────────────────────────────────────────────────────────
\titleformat{\section}[block]
  {\normalfont\large\bfseries}{}
  {0pt}{}[\vspace{2pt}\hrule\vspace{6pt}]
\titlespacing*{\section}{0pt}{14pt}{4pt}

% ── List formatting ─────────────────────────────────────────────────────────────
\setlist[enumerate,1]{label=\textbf{\arabic*.}, leftmargin=2em, itemsep=10pt}

% ── Branding constants (edit here, not per-file) ────────────────────────────────
\newcommand{\SpeedExamLine}{\textbf{Speed Maths:} Competition \& Exam Prep}
\newcommand{\SpeedCredit}{Series by \href{https://www.linkedin.com/in/cerealdev/}{David Akinyele-Aje}}

% ── Header / footer ──────────────────────────────────────────────────────────────
% Usage (once, after \input{../../shared/preamble} and before \begin{document}):
%   \SpeedHeader{Algebra}{3}
% First arg  = pillar name shown as "Speed <Pillar> --- Daily Drill"
% Second arg = worksheet number
\pagestyle{fancy}
\newcommand{\SpeedHeader}[2]{%
  \fancyhf{}%
  \renewcommand{\headrulewidth}{0.4pt}%
  \setlength{\headheight}{14pt}%
  \fancyhead[L]{\small\textbf{Speed #1 --- Daily Drill}}%
  \fancyhead[R]{\small Worksheet \##2}%
  \fancyfoot[L]{\small \SpeedExamLine}%
  \fancyfoot[C]{\small\thepage}%
  \fancyfoot[R]{\small \SpeedCredit}%
  \renewcommand{\footrulewidth}{0.4pt}%
}

% ── Title block (used at the top of every sheet AND answer file) ───────────────
% Usage: \SpeedTitleBlock{Daily Algebra Drill \#3}{Jane Doe}
% %% AGENTS: When generating a new sheet, ask the human contributor for the name they want credited in argument 2.
% For answer files, pass the "--- Answers \& Investigations" suffix in the arg.
\newcommand{\SpeedTitleBlock}[2]{%
  \begin{center}
    {\LARGE\bfseries #1}\\[4pt]
    {\normalsize\itshape Sheet by #2}\\[6pt]
    \hrule\vspace{4pt}
    \begin{tabular}{llcc}
      \toprule
      \textbf{Section} & \textbf{Topic} & \textbf{Questions} & \textbf{Target time}\\
      \midrule
      A & Rapid Recognition          & 10 & 2 min 30 sec \\
      B & Manipulation Drills        & 10 & 8 min        \\
      C & Substitution \& Structure  & 8  & 10 min       \\
      D & Challenge                  & 5  & 15 min       \\
      \bottomrule
    \end{tabular}
    \vspace{4pt}
    \hrule
  \end{center}
}

% ── Sheet metadata: topic + the day's new tools ────────────────────────────────
% Usage (immediately after \SpeedTitleBlock, sheets only):
%   \SpeedMeta{Expanding, factorising, and the identity toolkit}
%             {difference of squares, sum/product form, completing the square}
%
% Arg 1 = the sheet's topic in one line. The website lists this instead of
%         "Sheet 03", so write the thing a reader would actually search for.
% Arg 2 = comma-separated, at most five: the tools this sheet introduces that
%         earlier sheets in the pillar did not. These become the website's tags.
%         Leave out anything the reader already met — the tags are meant to say
%         what is new today, not to summarise the sheet.
%
% Both are parsed straight out of this file by tools/build_website.py, so the
% sheet stays the single source of truth and the site cannot drift from it.
%
% This renders NOTHING. It is a declaration for the build script, not a piece of
% the sheet's layout: adding it to a sheet must not change that sheet's PDF, so
% the 25 existing sheets can be annotated without a single one of them being
% reissued. If the toolkit should also appear on the page, that is a separate
% decision about the sheet design — make it deliberately, here, not by leaving
% this macro printing something nobody asked it to.
\newcommand{\SpeedMeta}[2]{}

% ── Answer / Method / Investigate-further boxes (answer files only) ────────────
\newmdenv[
  backgroundcolor=lightgray,
  linecolor=medgray,
  linewidth=0.5pt,
  leftmargin=0pt, rightmargin=0pt,
  innerleftmargin=8pt, innerrightmargin=8pt,
  innertopmargin=5pt, innerbottommargin=5pt,
  skipabove=4pt, skipbelow=4pt
]{ansbox}

\newmdenv[
  backgroundcolor=white,
  linecolor=darkgray,
  linewidth=0.8pt,
  leftline=true, rightline=false, topline=false, bottomline=false,
  leftmargin=0pt, rightmargin=0pt,
  innerleftmargin=8pt, innerrightmargin=6pt,
  innertopmargin=4pt, innerbottommargin=4pt,
  skipabove=4pt, skipbelow=6pt
]{invbox}

\newcommand{\ans}[1]{%
  \begin{ansbox}
    \textbf{Answer:} #1
  \end{ansbox}}

\newcommand{\method}[1]{%
  {\small\textbf{Method:} #1}\par}

\newcommand{\inv}[1]{%
  \begin{invbox}
    \textbf{\small Investigate further:} {\small #1}
  \end{invbox}}

% ── Misc helpers ────────────────────────────────────────────────────────────────
\newcommand{\qn}[1]{\textbf{#1.}\;}

% Mistake-linking: attach a Top 5 Patterns / Common Traps bullet to the question
% label(s) in this sheet where it actually shows up, e.g. \seealso{B6, D2}
\newcommand{\seealso}[1]{\hfill\textit{\footnotesize(see #1)}}

% ── Graph options macro ─────────────────────────────────────────────────────────
% Usage: \GraphOptions{tikz code for A}{tikz code for B}{tikz code for C}{tikz code for D}
\newcommand{\GraphOptions}[4]{%
  \begin{center}
    \begin{tabular}{cc}
      \textbf{A)} & \textbf{B)} \\
      #1 & #2 \\[10pt]
      \textbf{C)} & \textbf{D)} \\
      #3 & #4
    \end{tabular}
  \end{center}
}

% ── Closing motivational rule + cross-promo footer (sheets only) ────────────────
% Usage: \SpeedClosing{"Quote text."}
\newcommand{\SpeedClosing}[1]{%
  \vspace{20pt}
  \begin{center}
  \rule{0.6\textwidth}{0.4pt}\\[4pt]
  {\small\itshape #1}\\[4pt]
  \rule{0.6\textwidth}{0.4pt}
  \end{center}
  \begin{center}
  {\small Found a mistake or want to contribute a sheet?}\\
  {\small Visit the open-source project at \href{https://github.com/The-CerealDev/speed-maths}{github.com/The-CerealDev/speed-maths}}
  \end{center}
}
 and before \begin{document}):
%   \SpeedHeader{Algebra}{3}
% First arg  = pillar name shown as "Speed <Pillar> --- Daily Drill"
% Second arg = worksheet number
\pagestyle{fancy}
\newcommand{\SpeedHeader}[2]{%
  \fancyhf{}%
  \renewcommand{\headrulewidth}{0.4pt}%
  \setlength{\headheight}{14pt}%
  \fancyhead[L]{\small\textbf{Speed #1 --- Daily Drill}}%
  \fancyhead[R]{\small Worksheet \##2}%
  \fancyfoot[L]{\small \SpeedExamLine}%
  \fancyfoot[C]{\small\thepage}%
  \fancyfoot[R]{\small \SpeedCredit}%
  \renewcommand{\footrulewidth}{0.4pt}%
}

% ── Title block (used at the top of every sheet AND answer file) ───────────────
% Usage: \SpeedTitleBlock{Daily Algebra Drill \#3}{Jane Doe}
% %% AGENTS: When generating a new sheet, ask the human contributor for the name they want credited in argument 2.
% For answer files, pass the "--- Answers \& Investigations" suffix in the arg.
\newcommand{\SpeedTitleBlock}[2]{%
  \begin{center}
    {\LARGE\bfseries #1}\\[4pt]
    {\normalsize\itshape Sheet by #2}\\[6pt]
    \hrule\vspace{4pt}
    \begin{tabular}{llcc}
      \toprule
      \textbf{Section} & \textbf{Topic} & \textbf{Questions} & \textbf{Target time}\\
      \midrule
      A & Rapid Recognition          & 10 & 2 min 30 sec \\
      B & Manipulation Drills        & 10 & 8 min        \\
      C & Substitution \& Structure  & 8  & 10 min       \\
      D & Challenge                  & 5  & 15 min       \\
      \bottomrule
    \end{tabular}
    \vspace{4pt}
    \hrule
  \end{center}
}

% ── Sheet metadata: topic + the day's new tools ────────────────────────────────
% Usage (immediately after \SpeedTitleBlock, sheets only):
%   \SpeedMeta{Expanding, factorising, and the identity toolkit}
%             {difference of squares, sum/product form, completing the square}
%
% Arg 1 = the sheet's topic in one line. The website lists this instead of
%         "Sheet 03", so write the thing a reader would actually search for.
% Arg 2 = comma-separated, at most five: the tools this sheet introduces that
%         earlier sheets in the pillar did not. These become the website's tags.
%         Leave out anything the reader already met — the tags are meant to say
%         what is new today, not to summarise the sheet.
%
% Both are parsed straight out of this file by tools/build_website.py, so the
% sheet stays the single source of truth and the site cannot drift from it.
%
% This renders NOTHING. It is a declaration for the build script, not a piece of
% the sheet's layout: adding it to a sheet must not change that sheet's PDF, so
% the 25 existing sheets can be annotated without a single one of them being
% reissued. If the toolkit should also appear on the page, that is a separate
% decision about the sheet design — make it deliberately, here, not by leaving
% this macro printing something nobody asked it to.
\newcommand{\SpeedMeta}[2]{}

% ── Answer / Method / Investigate-further boxes (answer files only) ────────────
\newmdenv[
  backgroundcolor=lightgray,
  linecolor=medgray,
  linewidth=0.5pt,
  leftmargin=0pt, rightmargin=0pt,
  innerleftmargin=8pt, innerrightmargin=8pt,
  innertopmargin=5pt, innerbottommargin=5pt,
  skipabove=4pt, skipbelow=4pt
]{ansbox}

\newmdenv[
  backgroundcolor=white,
  linecolor=darkgray,
  linewidth=0.8pt,
  leftline=true, rightline=false, topline=false, bottomline=false,
  leftmargin=0pt, rightmargin=0pt,
  innerleftmargin=8pt, innerrightmargin=6pt,
  innertopmargin=4pt, innerbottommargin=4pt,
  skipabove=4pt, skipbelow=6pt
]{invbox}

\newcommand{\ans}[1]{%
  \begin{ansbox}
    \textbf{Answer:} #1
  \end{ansbox}}

\newcommand{\method}[1]{%
  {\small\textbf{Method:} #1}\par}

\newcommand{\inv}[1]{%
  \begin{invbox}
    \textbf{\small Investigate further:} {\small #1}
  \end{invbox}}

% ── Misc helpers ────────────────────────────────────────────────────────────────
\newcommand{\qn}[1]{\textbf{#1.}\;}

% Mistake-linking: attach a Top 5 Patterns / Common Traps bullet to the question
% label(s) in this sheet where it actually shows up, e.g. \seealso{B6, D2}
\newcommand{\seealso}[1]{\hfill\textit{\footnotesize(see #1)}}

% ── Graph options macro ─────────────────────────────────────────────────────────
% Usage: \GraphOptions{tikz code for A}{tikz code for B}{tikz code for C}{tikz code for D}
\newcommand{\GraphOptions}[4]{%
  \begin{center}
    \begin{tabular}{cc}
      \textbf{A)} & \textbf{B)} \\
      #1 & #2 \\[10pt]
      \textbf{C)} & \textbf{D)} \\
      #3 & #4
    \end{tabular}
  \end{center}
}

% ── Closing motivational rule + cross-promo footer (sheets only) ────────────────
% Usage: \SpeedClosing{"Quote text."}
\newcommand{\SpeedClosing}[1]{%
  \vspace{20pt}
  \begin{center}
  \rule{0.6\textwidth}{0.4pt}\\[4pt]
  {\small\itshape #1}\\[4pt]
  \rule{0.6\textwidth}{0.4pt}
  \end{center}
  \begin{center}
  {\small Found a mistake or want to contribute a sheet?}\\
  {\small Visit the open-source project at \href{https://github.com/The-CerealDev/speed-maths}{github.com/The-CerealDev/speed-maths}}
  \end{center}
}
 and before \begin{document}):
%   \SpeedHeader{Algebra}{3}
% First arg  = pillar name shown as "Speed <Pillar> --- Daily Drill"
% Second arg = worksheet number
\pagestyle{fancy}
\newcommand{\SpeedHeader}[2]{%
  \fancyhf{}%
  \renewcommand{\headrulewidth}{0.4pt}%
  \setlength{\headheight}{14pt}%
  \fancyhead[L]{\small\textbf{Speed #1 --- Daily Drill}}%
  \fancyhead[R]{\small Worksheet \##2}%
  \fancyfoot[L]{\small \SpeedExamLine}%
  \fancyfoot[C]{\small\thepage}%
  \fancyfoot[R]{\small \SpeedCredit}%
  \renewcommand{\footrulewidth}{0.4pt}%
}

% ── Title block (used at the top of every sheet AND answer file) ───────────────
% Usage: \SpeedTitleBlock{Daily Algebra Drill \#3}{Jane Doe}
% %% AGENTS: When generating a new sheet, ask the human contributor for the name they want credited in argument 2.
% For answer files, pass the "--- Answers \& Investigations" suffix in the arg.
\newcommand{\SpeedTitleBlock}[2]{%
  \begin{center}
    {\LARGE\bfseries #1}\\[4pt]
    {\normalsize\itshape Sheet by #2}\\[6pt]
    \hrule\vspace{4pt}
    \begin{tabular}{llcc}
      \toprule
      \textbf{Section} & \textbf{Topic} & \textbf{Questions} & \textbf{Target time}\\
      \midrule
      A & Rapid Recognition          & 10 & 2 min 30 sec \\
      B & Manipulation Drills        & 10 & 8 min        \\
      C & Substitution \& Structure  & 8  & 10 min       \\
      D & Challenge                  & 5  & 15 min       \\
      \bottomrule
    \end{tabular}
    \vspace{4pt}
    \hrule
  \end{center}
}

% ── Sheet metadata: topic + the day's new tools ────────────────────────────────
% Usage (immediately after \SpeedTitleBlock, sheets only):
%   \SpeedMeta{Expanding, factorising, and the identity toolkit}
%             {difference of squares, sum/product form, completing the square}
%
% Arg 1 = the sheet's topic in one line. The website lists this instead of
%         "Sheet 03", so write the thing a reader would actually search for.
% Arg 2 = comma-separated, at most five: the tools this sheet introduces that
%         earlier sheets in the pillar did not. These become the website's tags.
%         Leave out anything the reader already met — the tags are meant to say
%         what is new today, not to summarise the sheet.
%
% Both are parsed straight out of this file by tools/build_website.py, so the
% sheet stays the single source of truth and the site cannot drift from it.
%
% This renders NOTHING. It is a declaration for the build script, not a piece of
% the sheet's layout: adding it to a sheet must not change that sheet's PDF, so
% the 25 existing sheets can be annotated without a single one of them being
% reissued. If the toolkit should also appear on the page, that is a separate
% decision about the sheet design — make it deliberately, here, not by leaving
% this macro printing something nobody asked it to.
\newcommand{\SpeedMeta}[2]{}

% ── Answer / Method / Investigate-further boxes (answer files only) ────────────
\newmdenv[
  backgroundcolor=lightgray,
  linecolor=medgray,
  linewidth=0.5pt,
  leftmargin=0pt, rightmargin=0pt,
  innerleftmargin=8pt, innerrightmargin=8pt,
  innertopmargin=5pt, innerbottommargin=5pt,
  skipabove=4pt, skipbelow=4pt
]{ansbox}

\newmdenv[
  backgroundcolor=white,
  linecolor=darkgray,
  linewidth=0.8pt,
  leftline=true, rightline=false, topline=false, bottomline=false,
  leftmargin=0pt, rightmargin=0pt,
  innerleftmargin=8pt, innerrightmargin=6pt,
  innertopmargin=4pt, innerbottommargin=4pt,
  skipabove=4pt, skipbelow=6pt
]{invbox}

\newcommand{\ans}[1]{%
  \begin{ansbox}
    \textbf{Answer:} #1
  \end{ansbox}}

\newcommand{\method}[1]{%
  {\small\textbf{Method:} #1}\par}

\newcommand{\inv}[1]{%
  \begin{invbox}
    \textbf{\small Investigate further:} {\small #1}
  \end{invbox}}

% ── Misc helpers ────────────────────────────────────────────────────────────────
\newcommand{\qn}[1]{\textbf{#1.}\;}

% Mistake-linking: attach a Top 5 Patterns / Common Traps bullet to the question
% label(s) in this sheet where it actually shows up, e.g. \seealso{B6, D2}
\newcommand{\seealso}[1]{\hfill\textit{\footnotesize(see #1)}}

% ── Graph options macro ─────────────────────────────────────────────────────────
% Usage: \GraphOptions{tikz code for A}{tikz code for B}{tikz code for C}{tikz code for D}
\newcommand{\GraphOptions}[4]{%
  \begin{center}
    \begin{tabular}{cc}
      \textbf{A)} & \textbf{B)} \\
      #1 & #2 \\[10pt]
      \textbf{C)} & \textbf{D)} \\
      #3 & #4
    \end{tabular}
  \end{center}
}

% ── Closing motivational rule + cross-promo footer (sheets only) ────────────────
% Usage: \SpeedClosing{"Quote text."}
\newcommand{\SpeedClosing}[1]{%
  \vspace{20pt}
  \begin{center}
  \rule{0.6\textwidth}{0.4pt}\\[4pt]
  {\small\itshape #1}\\[4pt]
  \rule{0.6\textwidth}{0.4pt}
  \end{center}
  \begin{center}
  {\small Found a mistake or want to contribute a sheet?}\\
  {\small Visit the open-source project at \href{https://github.com/The-CerealDev/speed-maths}{github.com/The-CerealDev/speed-maths}}
  \end{center}
}

\SpeedHeader{Logic}{4}

\begin{document}

\SpeedTitleBlock{Daily Logic Drill \#4 --- Answers \& Investigations}{David Akinyele-Aje}
\noindent\textit{New toolkit today: Disproof by Counterexample $\cdot$ Boundary Case Discovery $\cdot$ Derivative \& Function Counterexamples $\cdot$ Universal Refutation.}


\section*{Section A}

\begin{enumerate}

%── A1 ──────────────────────────────────────────────────────────────────────────
\item Provide the minimal positive integer counterexample $n$ to: ``For all positive integers $n$, if $n$ is prime, then $n^2+2$ is not prime.'' \textit{\small(after TMUA 2022 Paper 2 Q3)}

\ans{$3$}
\method{Test small primes: For $n=2$, $2^2+2=6$ (composite, not a counterexample). For $n=3$, $3$ is prime and $3^2+2 = 11$, which is prime! Thus $n=3$ satisfies the premise but falsifies the conclusion, making it the minimal counterexample.}
\inv{For any prime $p > 3$, $p \equiv \pm 1 \pmod 3$, so $p^2 \equiv 1 \pmod 3$, which means $p^2+2 \equiv 0 \pmod 3$, so $p^2+2$ is always a multiple of $3$ and thus composite. Hence $n=3$ is the \emph{only} prime counterexample!}

%── A2 ──────────────────────────────────────────────────────────────────────────
\item Find the smallest positive integer $a$ for which there exist positive integers $b, c$ such that $a \mid bc$ but $a \nmid b$ and $a \nmid c$. \textit{\small(after TMUA 2021 Paper 2 Q4)}

\ans{$4$}
\method{If $a$ is prime, Euclid's Lemma guarantees $a \mid bc \implies a \mid b \lor a \mid c$. Thus $a$ must be composite. The smallest composite number is $a=4$. Taking $b=2$ and $c=2$, we have $bc = 4$, so $4 \mid 4$, but $4 \nmid 2$. Thus $a=4$ is the minimal counterexample.}
\inv{If $a$ must be squarefree, what is the smallest counterexample? $a=6$, with $b=2, c=3$.}

%── A3 ──────────────────────────────────────────────────────────────────────────
\item Provide a function $f(x)$ with $f'(x) > 0$ for all $x \in \mathbb{R}$ that serves as a counterexample to: ``If $f'(x) > 0$ for all $x$, then $f(x) > 0$ for all $x$.'' \textit{\small(after TMUA 2018 Paper 2 Q6)}

\ans{$f(x) = x$}
\method{The linear function $f(x) = x$ has $f'(x) = 1 > 0$ for all $x \in \mathbb{R}$, but $f(x) \le 0$ for all $x \le 0$ (e.g.\ $f(-5) = -5 \not> 0$). Another valid counterexample is $f(x) = e^x - 2$.}
\inv{A positive derivative guarantees that the function is strictly increasing, not that its values are positive.}

%── A4 ──────────────────────────────────────────────────────────────────────────
\item Find a real number $x > 0$ that serves as a counterexample to: ``$x > \frac{1}{x}$ for all $x > 0$.''

\ans{$\frac{1}{2}$}
\method{For $x = \frac{1}{2} > 0$, we have $\frac{1}{x} = 2$. Since $\frac{1}{2} < 2$, the inequality fails. Any $x \in (0, 1]$ disproves the claim.}
\inv{The inequality $x > 1/x$ on positive reals holds if and only if $x > 1$.}

%── A5 ──────────────────────────────────────────────────────────────────────────
\item Provide a pair of real numbers $(x, y)$ that serves as a counterexample to: ``For all $x, y \in \mathbb{R}$, if $x^2 > y^2$, then $x > y$.''

\ans{$(-3, 1)$}
\method{Let $x = -3$ and $y = 1$. Then $x^2 = 9$ and $y^2 = 1$, so $x^2 > y^2$ holds. However, $-3 < 1$, so $x > y$ is false.}
\inv{Squaring preserves inequality only when both numbers are non-negative.}

%── A6 ──────────────────────────────────────────────────────────────────────────
\item Find the smallest positive integer $n$ for which $2^n - 1$ is composite despite $n$ being prime.

\ans{$11$}
\method{Primes $n$: $2^2-1=3$ (prime), $2^3-1=7$ (prime), $2^5-1=31$ (prime), $2^7-1=127$ (prime). For $n=11$: $2^{11}-1 = 2047 = 23 \times 89$, which is composite. Thus $n=11$ is the smallest counterexample.}
\inv{Mersenne numbers $2^p-1$ need not be prime even when $p$ is prime.}

%── A7 ──────────────────────────────────────────────────────────────────────────
\item Provide a pair of non-negative real numbers $(a, b)$ that serves as a counterexample to: ``$\sqrt{a+b} = \sqrt{a} + \sqrt{b}$.''

\ans{$(1, 1)$}
\method{For $a=1, b=1$, $\sqrt{a+b} = \sqrt{2}$, whereas $\sqrt{a}+\sqrt{b} = 1+1=2$. Since $\sqrt{2} \neq 2$, the identity fails.}
\inv{The equality $\sqrt{a+b}=\sqrt{a}+\sqrt{b}$ holds if and only if $ab = 0$.}

%── A8 ──────────────────────────────────────────────────────────────────────────
\item Find the smallest positive integer $n$ that serves as a counterexample to: ``$n^2 + 1$ is always prime.''

\ans{$3$}
\method{For $n=1$, $1^2+1=2$ (prime). For $n=2$, $2^2+1=5$ (prime). For $n=3$, $3^2+1=10 = 2 \times 5$, which is composite. Smallest $n$ is $3$.}
\inv{For which other small integers is $n^2+1$ composite? $n=5 \implies 26 = 2 \times 13$, $n=7 \implies 50$.}

%── A9 ──────────────────────────────────────────────────────────────────────────
\item Provide a differentiable function $f: [1, \infty) \to \mathbb{R}$ with $\lim_{x\to\infty} f'(x) = 0$ that disproves: ``If $\lim_{x\to\infty} f'(x) = 0$, then $\lim_{x\to\infty} f(x)$ exists and is finite.''

\ans{$f(x) = \sqrt{x}$}
\method{For $f(x) = \sqrt{x}$, $f'(x) = \frac{1}{2\sqrt{x}} \to 0$ as $x \to \infty$, but $f(x) = \sqrt{x} \to \infty$. Another standard counterexample is $f(x) = \ln x$, where $f'(x) = 1/x \to 0$ while $\ln x \to \infty$.}
\inv{A vanishing derivative means the growth rate slows to zero, but the total accumulated growth can still diverge to infinity.}

%── A10 ─────────────────────────────────────────────────────────────────────────
\item Find the smallest odd composite integer $n$ that is not a prime power.

\ans{$15$}
\method{Odd composites in order: $9 = 3^2$ (prime power), $15 = 3 \times 5$ (not a prime power). Thus $15$ is the smallest.}
\inv{Any non-prime-power odd composite has at least two distinct odd prime factors.}

\end{enumerate}

\section*{Section B}

\begin{enumerate}

%── B1 ──────────────────────────────────────────────────────────────────────────
\item A student claims: ``If a polynomial $f(x)$ satisfies $f(n) \in \mathbb{Z}$ for all $n \in \mathbb{Z}$, then $f'(n) \in \mathbb{Z}$ for all $n \in \mathbb{Z}$.'' Provide a quadratic polynomial $f(x)$ that serves as a counterexample. \textit{\small(after TMUA 2017 Paper 2 Q16)}

\ans{$f(x) = \frac{x(x-1)}{2}$}
\method{For $f(x) = \frac{x^2-x}{2} = \binom{x}{2}$, the product of two consecutive integers $n(n-1)$ is always even, so $f(n) \in \mathbb{Z}$ for all $n \in \mathbb{Z}$. However, $f'(x) = x - \frac{1}{2}$. For any integer $n$, $f'(n) = n - \frac{1}{2} \notin \mathbb{Z}$. This disproves the claim.}
\inv{This counterexample relies on the binomial basis polynomials $\binom{x}{k}$, which are integer-valued on $\mathbb{Z}$ but do not have integer coefficients.}

%── B2 ──────────────────────────────────────────────────────────────────────────
\item Give a pair of irrational numbers $a, b$ whose sum is rational.

\ans{$a = \sqrt{2}, b = -\sqrt{2}$}
\method{$\sqrt{2}$ and $-\sqrt{2}$ are both irrational, but their sum is $\sqrt{2} + (-\sqrt{2}) = 0 \in \mathbb{Q}$.}
\inv{Another pair: $a = 1 + \sqrt{3}, b = 1 - \sqrt{3}$, where $a+b = 2$.}

%── B3 ──────────────────────────────────────────────────────────────────────────
\item Give a pair of irrational numbers $a, b$ whose product is rational.

\ans{$a = \sqrt{2}, b = \sqrt{2}$}
\method{$\sqrt{2}$ is irrational, but $ab = \sqrt{2} \times \sqrt{2} = 2 \in \mathbb{Q}$.}
\inv{A pair of distinct positive irrationals: $a = \sqrt{8}, b = \sqrt{2}$, where $ab = \sqrt{16} = 4$.}

%── B4 ──────────────────────────────────────────────────────────────────────────
\item Provide a pair of functions $f, g: [0, \infty) \to \mathbb{R}$ such that $f(x) \ge g(x)$ for all $x \ge 0$, but $f'(0) < g'(0)$. \textit{\small(after TMUA 2021 Paper 2 Q12)}

\ans{$f(x) = 10, g(x) = x$}
\method{For $x \in [0, 5]$, $f(x) = 10 \ge x = g(x)$. To make $f(x) \ge g(x)$ for all $x \ge 0$, let $f(x) = 10$ and $g(x) = 10 - e^{-x}$. Then $f(x) = 10 > g(x)$ for all $x \ge 0$, but $f'(0) = 0$ while $g'(0) = 1 > 0$, so $f'(0) < g'(0)$.}
\inv{Inequalities between functions cannot be differentiated; a function can sit entirely above another while having a smaller derivative.}

%── B5 ──────────────────────────────────────────────────────────────────────────
\item Find the smallest positive integer $n$ that is a counterexample to: ``$n^2 + n + 41$ is prime for all positive integers $n$.''

\ans{$40$}
\method{For $n=40$, $n^2+n+41 = 40^2+40+41 = 40(41)+41 = 41(40+1) = 41^2 = 1681$, which is composite ($41 \times 41$). For all $1 \le n \le 39$, $n^2+n+41$ is prime. Thus $n=40$ is the minimal counterexample.}
\inv{At $n=41$, $41^2+41+41 = 41 \times 43$, which is also composite, but $40$ is smaller.}

%── B6 ──────────────────────────────────────────────────────────────────────────
\item Provide a counterexample of three finite sets $A, B, C$ showing that $A \cup B = A \cup C$ does not imply $B = C$.

\ans{$A = \{1\}, B = \{1\}, C = \emptyset$}
\method{Here $A \cup B = \{1\} \cup \{1\} = \{1\}$, and $A \cup C = \{1\} \cup \emptyset = \{1\}$. The unions are equal, but $B = \{1\} \neq \emptyset = C$.}
\inv{Union is not an invertible operation: elements absorbed by $A$ lose their distinguishing presence in $B$ or $C$.}

%── B7 ──────────────────────────────────────────────────────────────────────────
\item Find the smallest positive integer $n$ for which $3^n + 2$ is composite.

\ans{$5$}
\method{Test values of $n$:
$n=1 \implies 3^1+2 = 5$ (prime);
$n=2 \implies 3^2+2 = 11$ (prime);
$n=3 \implies 3^3+2 = 29$ (prime);
$n=4 \implies 3^4+2 = 83$ (prime);
$n=5 \implies 3^5+2 = 243+2 = 245 = 5 \times 49$ (composite).
The smallest positive integer is $n=5$.}
\inv{Modulo $5$: $3^1 \equiv 3, 3^2 \equiv 4, 3^3 \equiv 2, 3^4 \equiv 1$, so $3^5 \equiv 3 \implies 3^5+2 \equiv 5 \equiv 0 \pmod 5$.}

%── B8 ──────────────────────────────────────────────────────────────────────────
\item A student conjectures: ``For every odd prime $p$, the integer $2^p - p$ is divisible by either $3$ or $5$.'' Find the smallest odd prime $p$ that serves as a counterexample. \textit{\small(after TMUA 2021 Paper 2 Q10)}

\ans{$7$}
\method{Test odd primes:
\begin{itemize}
\item $p=3$: $2^3 - 3 = 8 - 3 = 5$ (divisible by $5$).
\item $p=5$: $2^5 - 5 = 32 - 5 = 27$ (divisible by $3$).
\item $p=7$: $2^7 - 7 = 128 - 7 = 121 = 11^2$. $121$ is divisible by neither $3$ nor $5$.
\end{itemize}
Thus $p=7$ is the smallest odd prime counterexample.}
\inv{Small patterns in modular arithmetic often break precisely at $p=7$, the first prime where neither $2^p \equiv p \pmod 3$ nor $2^p \equiv p \pmod 5$ holds.}

%── B9 ──────────────────────────────────────────────────────────────────────────
\item Provide integers $a, b$ and a modulus $m > 1$ showing that $a^2 \equiv b^2 \pmod m$ does not imply $a \equiv \pm b \pmod m$.

\ans{$m=8, a=1, b=3$}
\method{For $m=8$, $a=1, b=3$: $a^2 = 1 \equiv 1 \pmod 8$ and $b^2 = 9 \equiv 1 \pmod 8$, so $a^2 \equiv b^2 \pmod 8$. However, $3 \not\equiv 1 \pmod 8$ and $3 \not\equiv -1 \equiv 7 \pmod 8$.}
\inv{The failure occurs because $8$ is composite; in $\mathbb{Z}/8\mathbb{Z}$, the polynomial $x^2-1$ has four roots $\{1, 3, 5, 7\}$.}

%── B10 ─────────────────────────────────────────────────────────────────────────
\item Give a counterexample function $f: [0, 1] \to [0, 1]$ showing that a continuous function need not have a UNIQUE fixed point.

\ans{$f(x) = x$}
\method{The identity function $f(x) = x$ maps $[0, 1]$ continuously to $[0, 1]$. Every single $x \in [0, 1]$ is a fixed point, meaning there are infinitely many fixed points rather than a unique one.}
\inv{Brouwer's Fixed Point Theorem guarantees the \emph{existence} of at least one fixed point, but never uniqueness.}

\end{enumerate}

\section*{Section C}

\begin{enumerate}

%── C1 ──────────────────────────────────────────────────────────────────────────
\item Consider the statement: ``$(*)$ If a whole number $n$ is $2$ more or $4$ more than a multiple of $5$, then $n$ is prime.'' How many counterexamples to $(*)$ are there in the range $1\leq n\leq50$? \textit{\small(after TMUA 2019 Paper 2 Q4)}

\ans{$13$}
\method{Numbers of the form $5k+2$ in $[1, 50]$ are: $2, 7, 12, 17, 22, 27, 32, 37, 42, 47$. The composites among these are $12, 22, 27, 32, 42$ ($5$ counterexamples).\\
Numbers of the form $5k+4$ in $[1, 50]$ are: $4, 9, 14, 19, 24, 29, 34, 39, 44, 49$. The composites among these are $4, 9, 14, 24, 34, 39, 44, 49$ ($8$ counterexamples).\\
Total counterexamples $= 5 + 8 = 13$.}
\inv{Note that $2$ and $7$ satisfy the premise and are prime, so they are not counterexamples.}

%── C2 ──────────────────────────────────────────────────────────────────────────
\item Which of the following values of $n$ is a counterexample to the claim: ``For all positive integers $n$, if $n$ is prime, then $n^3 - n$ is NOT divisible by $24$?'' \textit{\small(after TMUA 2022 Paper 2 Q3)}

\ans{E) Both $n=5$ and $n=7$.}
\method{Factor $n^3-n = (n-1)n(n+1)$. For any prime $p \ge 5$, $p$ is odd, so $p-1$ and $p+1$ are consecutive even integers, one of which is divisible by $4$, so their product is divisible by $8$. Furthermore, among $p-1, p, p+1$, exactly one is divisible by $3$; since $p \ge 5$ is prime, $3 \nmid p$, so $3 \mid (p-1)(p+1)$. Thus for any prime $p \ge 5$, $24 \mid (p^3-p)$. Both $n=5$ ($5^3-5=120 = 24 \times 5$) and $n=7$ ($7^3-7=336 = 24 \times 14$) are divisible by $24$, refuting the claim. Option E.}
\inv{This means every single prime $p \ge 5$ is a counterexample to the claim!}

%── C3 ──────────────────────────────────────────────────────────────────────────
\item A student claims: ``If $f:\mathbb{R}\to\mathbb{R}$ is twice differentiable, $f(0)=0$, and $f''(x) > 0$ for all $x \in \mathbb{R}$, then $f(x) \ge 0$ for all $x \ge 0$.'' Which of the following functions serves as a counterexample? \textit{\small(after TMUA 2018 Paper 2 Q6)}

\ans{C) $f(x) = x^2 - 3x$.}
\method{Check the conditions for $f(x) = x^2 - 3x$:
$f(0) = 0$; $f'(x) = 2x - 3$; $f''(x) = 2 > 0$ for all $x$.
All hypotheses hold. However, for $x \in (0, 3)$, $f(x) = x(x-3) < 0$. For instance, $f(1) = -2 < 0$. This disproves the claim. Option C.}
\inv{A positive second derivative means the graph is convex (bending upward), but if the initial slope $f'(0)$ is negative, the function dips below zero before rising.}

%── C4 ──────────────────────────────────────────────────────────────────────────
\item Consider the statement: ``For all positive integers $a, b, c$, if $a \mid bc$, then $a \mid b$ or $a \mid c$.'' Which of the following triples $(a, b, c)$ is a valid counterexample? \textit{\small(after TMUA 2021 Paper 2 Q4)}

\ans{B) $(6, 4, 9)$.}
\method{Check option B: $bc = 4 \times 9 = 36$. $a = 6$ divides $36$ ($36/6 = 6$). However, $6 \nmid 4$ and $6 \nmid 9$. Both disjuncts fail while the premise holds. Thus $(6, 4, 9)$ is a valid counterexample. In option A, $a=2 \mid 4$; in C, $a=5 \mid 10$; in D, $a=7 \mid 14$.}
\inv{A counterexample requires $a$ to share distinct prime factors between $b$ and $c$.}

%── C5 ──────────────────────────────────────────────────────────────────────────
\item Which of the following polynomials satisfies $f(n) \in \mathbb{Z}$ for all $n \in \mathbb{Z}$, but $f'(0) \notin \mathbb{Z}$? \textit{\small(after TMUA 2017 Paper 2 Q16)}

\ans{D) Both B and C.}
\method{For B: $f(x) = \frac{x^2+x}{2} = \binom{x+1}{2}$, integer-valued on $\mathbb{Z}$. $f'(x) = x + \frac{1}{2} \implies f'(0) = \frac{1}{2} \notin \mathbb{Z}$.\\
For C: $f(x) = \frac{x^3-x}{6} = \frac{(x-1)x(x+1)}{6} = \binom{x+1}{3}$, integer-valued on $\mathbb{Z}$ since the product of three consecutive integers is always divisible by $6$. $f'(x) = \frac{3x^2-1}{6} \implies f'(0) = -\frac{1}{6} \notin \mathbb{Z}$.\\
Both B and C satisfy the condition. Option D.}
\inv{Every integer-valued polynomial is an integer linear combination of the binomial polynomials $\binom{x}{k}$.}

%── C6 ──────────────────────────────────────────────────────────────────────────
\item A student conjectures: ``If a continuous function $f:[-\pi, \pi] \to \mathbb{R}$ satisfies $\int_{-\pi}^\pi f(x)dx = 0$, then $f(x) = 0$ for all $x \in [-\pi, \pi]$.'' Which of the following is a valid counterexample? \textit{\small(after TMUA 2019 Paper 2 Q14)}

\ans{D) All of A, B, C.}
\method{Evaluate the integral for each function on $[-\pi, \pi]$:
\begin{itemize}
\item A: $\int_{-\pi}^\pi \cos x\,dx = [\sin x]_{-\pi}^\pi = 0 - 0 = 0$, but $\cos 0 = 1 \neq 0$.
\item B: $\int_{-\pi}^\pi \sin x\,dx = [-\cos x]_{-\pi}^\pi = 1 - 1 = 0$, but $\sin(\pi/2) = 1 \neq 0$.
\item C: $\int_{-\pi}^\pi x\,dx = \left[\frac{x^2}{2}\right]_{-\pi}^\pi = 0$, but $f(1) = 1 \neq 0$.
\end{itemize}
All three functions have zero integral yet are not identically zero. Option D.}
\inv{Any non-trivial odd function on a symmetric interval $[-a, a]$ integrates to zero.}

%── C7 ──────────────────────────────────────────────────────────────────────────
\item Consider two statements about functions $f, g: [0, \infty) \to \mathbb{R}$ with $f(0)=g(0)=0$:
\begin{itemize}
\item[I.] $f'(x) \ge g'(x) \text{ for all } x \ge 0 \implies f(x) \ge g(x) \text{ for all } x \ge 0$.
\item[II.] $f(x) \ge g(x) \text{ for all } x \ge 0 \implies f'(x) \ge g'(x) \text{ for all } x \ge 0$.
\end{itemize}
Which of I and II are true? \textit{\small(after TMUA 2021 Paper 2 Q12)}

\ans{B) I only.}
\method{By the Fundamental Theorem of Calculus, $f(x) - g(x) = \int_0^x (f'(t) - g'(t))\,dt$. If $f'(t) \ge g'(t)$, the integrand is non-negative, so the integral is $\ge 0$, proving $f(x) \ge g(x)$. Thus statement I is TRUE.\\
For statement II: consider $f(x) = 2x$ and $g(x) = x + \sin x$. For $x \ge 0$, $f(x) - g(x) = x - \sin x \ge 0$, so $f(x) \ge g(x)$. However, $f'(x) = 2$ and $g'(x) = 1 + \cos x$. At $x = 0$, $f'(0) = 2 \ge 2$, but consider $f(x) = 10x$ and $g(x) = 10x - 5\cos(10x) + 5$: oscillating derivatives frequently exceed $f'(x)$. Statement II is FALSE. Option B.}
\inv{Integration preserves inequalities; differentiation does not.}

%── C8 ──────────────────────────────────────────────────────────────────────────
\item Consider the theorem: ``An odd prime $p$ can be written as the sum of two integer squares if and only if $p \equiv 1 \pmod 4$.'' Which of the following would constitute a counterexample to the \textbf{converse} of this theorem? \textit{\small(after TMUA 2018 Paper 2 Q17)}

\ans{A) A prime $p \equiv 1 \pmod 4$ that cannot be written as the sum of two squares.}
\method{The theorem states $P \iff Q$, where $P$: ``$p$ is a sum of two squares'' and $Q$: ``$p \equiv 1 \pmod 4$'' (for odd primes). The converse is $Q \Rightarrow P$: ``If an odd prime $p \equiv 1 \pmod 4$, then $p$ can be written as the sum of two squares.'' A counterexample to $Q \Rightarrow P$ must satisfy the hypothesis $Q$ ($p \equiv 1 \pmod 4$) but fail the conclusion $P$ (cannot be written as the sum of two squares). Option A.}
\inv{Refuting a conditional statement requires satisfying the antecedent while falsifying the consequent.}

\end{enumerate}

\section*{Section D}

\begin{enumerate}

%── D1 ──────────────────────────────────────────────────────────────────────────
\item Consider: I. ``$n=41$ is a counterexample to `$n^2+n+41$ is always prime'.'' II. ``$n=40$ is the smallest counterexample to `$n^2+n+41$ is always prime'.'' III. ``$n=4$ is a counterexample to `$n^2+n+41$ is always prime'.'' Which of I, II, III are true?

\ans{E) I and II only}
\method{I: $41^2+41+41 = 41(41+1+1) = 41 \times 43$, composite. So I is true.\\
II: For $1 \le n \le 39$, $n^2+n+41$ is prime. At $n=40$, $40^2+40+41 = 41^2$, composite. Thus $n=40$ is indeed the smallest positive integer counterexample. II is true.\\
III: For $n=4$, $4^2+4+41 = 61$, which is prime. So $n=4$ is not a counterexample. III is false.\\
Hence, I and II only are true.}
\inv{Euler's polynomial $n^2+n+41$ is the classical demonstration that finite verification cannot prove universal claims.}

%── D2 ──────────────────────────────────────────────────────────────────────────
\item Prove by contradiction: ``There do not exist positive integers $a,b$ with $a^2-b^2=1$.''

\ans{Proof by contradiction}
\method{Assume toward a contradiction that there exist positive integers $a,b$ such that $a^2 - b^2 = 1$. Factorising as a difference of two squares gives $(a-b)(a+b) = 1$. Since $a,b \in \mathbb{Z}^+$, we have $a \ge 1$ and $b \ge 1$, so $a+b \ge 2$. However, the only factorisation of $1$ into positive integers is $1 \times 1$, requiring $a+b = 1$. This contradicts $a+b \ge 2$. Thus, no such positive integers exist.}
\inv{Do non-zero integer solutions exist if $a,b$ are not restricted to positive integers?}

%── D3 ──────────────────────────────────────────────────────────────────────────
\item A student claims: ``For every positive integer $n$, at least one of $n,n+2,n+4$ is divisible by $3$.'' Determine, with proof, whether this claim is true.

\ans{True}
\method{Proof: Any positive integer $n$ can be written as $3k$, $3k+1$, or $3k+2$ for some integer $k \ge 0$.\\
Case 1: If $n = 3k$, then $n$ is divisible by $3$.\\
Case 2: If $n = 3k+1$, then $n+2 = (3k+1)+2 = 3k+3 = 3(k+1)$, which is divisible by $3$.\\
Case 3: If $n = 3k+2$, then $n+4 = (3k+2)+4 = 3k+6 = 3(k+2)$, which is divisible by $3$.\\
In all cases, at least one of $n, n+2, n+4$ is divisible by $3$. The claim is true.}
\inv{Notice that $n, n+2, n+4 \equiv n, n-1, n+1 \pmod 3$, so they cover all three residue classes modulo $3$.}

%── D4 ──────────────────────────────────────────────────────────────────────────
\item Give a counterexample to: ``If $p$ is prime, then $2^p+1$ is prime.''

\ans{$3$}
\method{For $p=3$ (which is prime), $2^p+1 = 2^3+1 = 8+1 = 9 = 3^2$, which is composite. (Also for $p=5$, $2^5+1 = 33 = 3 \times 11$).}
\inv{For any odd prime $p$, $2^p+1 = (2+1)(2^{p-1}-2^{p-2}+\dots+1) = 3 \times M$, so $2^p+1$ is always a multiple of $3$, and hence composite for all $p > 1$ with $2^p+1 > 3$!}

%── D5 ──────────────────────────────────────────────────────────────────────────
\item Statement $(*)$: ``For every positive integer $n\geq2$, there is a prime strictly between $n$ and $2n$.'' (a) Verify $(*)$ for $n=2,3,4,5$ by direct computation. (b) Does this verification constitute a proof of $(*)$? Justify your answer.

\ans{(a) $n=2: 3$; $n=3: 5$; $n=4: 5$ or $7$; $n=5: 7$. (b) No, finite checking does not prove a universal statement.}
\method{(a) For $n=2$, prime $3 \in (2,4)$. For $n=3$, prime $5 \in (3,6)$. For $n=4$, prime $5 \in (4,8)$. For $n=5$, prime $7 \in (5,10)$. \\
(b) No. Verification for finitely many cases ($n=2,3,4,5$) does not constitute a proof for infinitely many positive integers $n \ge 2$. A single unchecked value could be a counterexample. (The statement is true for all $n \ge 2$, known as Bertrand's Postulate, but proving it requires a general argument).}
\inv{Bertrand's Postulate was first proved by Chebyshev in 1852 using properties of the gamma function and prime factorisations of central binomial coefficients.}

\end{enumerate}

%═══════════════════════════════════════════════════════════════════════════════
\section*{Top 5 Patterns Today}
%═══════════════════════════════════════════════════════════════════════════════

\begin{enumerate}[leftmargin=2em, itemsep=4pt]
  \item \textbf{Refuting Implications by Finding Witness Instances.}
        To disprove $\forall x, (P(x) \Rightarrow Q(x))$, find a single counterexample where $P(x)$ holds while $Q(x)$ fails. \seealso{A1, A2, A3, A4, B8}
  \item \textbf{Derivative and Function Inversion Traps.}
        A function satisfying $f'(x)>0$ everywhere does not need to be positive ($f(x)=x$); similarly, $f(x) \ge g(x)$ does not force $f'(x) \ge g'(x)$. \seealso{A3, A9, B4, C3, C7}
  \item \textbf{Integer-Valued Polynomial Derivatives.}
        A polynomial integer-valued on $\mathbb{Z}$ (such as binomial coefficients $\binom{x}{2} = \frac{x^2-x}{2}$) need not have integer coefficients or integer derivatives. \seealso{B1, C5}
  \item \textbf{Prime Generator Boundaries and Divisibility.}
        Conjectures asserting primality (like $n^2+n+41$ or $2^p-p$) fail at structural boundaries ($n=40$, $p=7$) where algebraic factoring or non-residues intervene. \seealso{A6, A8, B5, B7, D1}
  \item \textbf{Non-Trivial Modular Residue Square Roots.}
        In non-prime moduli (like $m=8$), the congruence $x^2 \equiv 1 \pmod 8$ has four distinct roots ($\pm 1, \pm 3$), breaking the field property $a^2 \equiv b^2 \implies a \equiv \pm b$. \seealso{A10, B9, C4}
\end{enumerate}

%═══════════════════════════════════════════════════════════════════════════════
\section*{Common Traps to Avoid}
%═══════════════════════════════════════════════════════════════════════════════

\begin{enumerate}[leftmargin=2em, itemsep=4pt]
  \item \textbf{Testing Only Positive or Symmetrical Values.}
        Order relations like $x^2>y^2 \implies x>y$ fail only when considering negative values ($x=-3, y=1$); always check signed edge cases. \seealso{A4, A5, A7, B2}
  \item \textbf{Confusing Implication Refutation with Converse Refutation.}
        Refuting the converse of $P \iff Q$ requires finding a witness satisfying $Q$ but not $P$, rather than attempting to falsify the forward direction. \seealso{C8, A1, B8}
  \item \textbf{Assuming Derivatives Inherit Function Inequalities.}
        The height of a curve $f(x) \ge g(x)$ provides zero information about its slope $f'(x) \ge g'(x)$; slopes can cross while functions remain separated. \seealso{B4, C3, C7}
  \item \textbf{Assuming Finite Verification Equals Proof.}
        Checking that a statement holds for hundreds of small cases (Bertrand's Postulate, Euler's polynomial) demonstrates only the absence of small counterexamples, never a proof. \seealso{D5, B5, D1}
  \item \textbf{Overlooking Composite Prime Factors $p \le \sqrt{n}$.}
        Refuting that $a \mid bc \implies a \mid b \lor a \mid c$ requires a composite modulus $a$ that splits its prime factors across $b$ and $c$. \seealso{A2, C4, B9}
\end{enumerate}

\SpeedClosing{``Speed comes from recognition, not from rushing. Every shortcut is a pattern you have internalised.''}

\end{document}
